\section{Integração e verificação de coerência}
O problema central gera atrasos, ruídos e dispersão que afetam emissores, avaliadores e destinatários. No processo AS-IS, a dor aparece na utilização de canais informais, na falta de moderação uniforme e na ausência de rastreamento de quem recebeu e visualizou a comunicação.

Os requisitos RF01 a RF06 reestruturam esse fluxo: o emissor cria avisos, o sistema verifica necessidade de autorização, o avaliador decide, o resultado é comunicado e o destinatário recebe a informação segmentada. As histórias e cenários BDD especificam fluxos principal e alternativo. Finalmente, os assuntos organizados dão suporte a esses comportamentos.

\begin{table}[!ht]\centering\small\renewcommand{\arraystretch}{1.45}
\begin{tabularx}{\linewidth}{|>{\raggedright\arraybackslash}X|>{\raggedright\arraybackslash}X|}
\hline
\textbf{Ponto de verificação} & \textbf{Situação no projeto} \\
\hline
Requisitos ligados às necessidades & Todos os RF foram derivados de stakeholders e etapas do BPMN. \\
\hline
Dados suficientes para os RFs & Aviso, usuário, avaliação, segmentação, envio e solicitação cobrem os fluxos. \\
\hline
Histórias e BDD coerentes & Cada requisito possui história e cenários de sucesso e exceção. \\
\hline
Processo futuro melhora o AS-IS & A centralização e a segmentação reduzem dispersão, perda e dependência de canais informais. \\
\hline
\end{tabularx}
\end{table}
